\documentclass[10pt, a4paper]{article}
\usepackage{preamble}

\title{Probability II}
\author{Luke Phillips}
\date{August 2025}

\begin{document}

\maketitle

\newpage

\tableofcontents

\newpage

\section{Probability and independence}

\subsection{Probability spaces}

\begin{definition}[Sigma-field]
    Let $\mathcal{F}$ be a collection of subsets of $\Omega$.
    We shall call $\mathcal{F}$ a $\sigma$-field or a $\sigma$-algebra if it has the following properties:
    \begin{enumerate}[label = \arabic*.]
        \item $\varnothing \in \mathcal{F}$

        \item if $A \in \mathcal{F}$,
        then $A ^ c \in \mathcal{F}$,

        \item if $A_1, A_2, \dotsc \in \mathcal{F}$,
        then $\bigcup_{k = 1}^{\infty}A_k \in \mathcal{F}$.
    \end{enumerate}
\end{definition}

\begin{definition}[Probability]
    Let $\Omega$ be a sample space,
    and $\mathcal{F}$ be a $\sigma$-field of events in $\Omega$.
    A probability distribution $\P$ on $(\Omega, \mathcal{F})$ is a collection of numbers $\P(A)$,
    $A \in \mathcal{F}$,
    possessing the following properties:
    \begin{enumerate}[label = A\arabic*]
        \item  for every event $A \in \mathcal{F}$,
        $\P(A) \geq 0$

        \item $\P(\Omega) = 1$

        \item for any pair of disjoint events $A$ and $B$
        (i.e.,
        $A \cap B = \emptyset$),
        $\P(A \cup B) = \P(A) + \P(B)$

        \item for any countable collection of $A_1, A_2, \dotsc$ of disjoint events
        (i.e.,
        with $A_i \cap A_j = \varnothing$ for $i \neq j$),
        \[
        \P\left(\bigcup_{k = 1}^{\infty}A_k\right) = \sum_{k = 1}^{\infty}\P(A_k).
        \]
    \end{enumerate}
\end{definition}

\begin{definition}
    A probability space is a triple $(\Omega, \mathcal{F}, \P)$,
    where $\Omega$ is a sample space,
    $\mathcal{F}$ is a $\sigma$-field of events in $\Omega$,
    and $\P(\cdot)$ is a probability measure on $(\Omega, \mathcal{F})$.
\end{definition}

\subsubsection{Generated \texorpdfstring{$\sigma$}{}-fields}
\begin{definition}\label{prob:def:siggenfldD}
    Let $\mathcal{X}$ be an arbitrary set and let $\mathcal{D}$ be an arbitrary collection of subsets of $\mathcal{X}$.
    Let $\mathcal{F}_{\alpha}, \alpha \in \mathcal{A}$,
    be the collection of all sigma-fields in $\mathcal{X}$ which contain $\mathcal{D}$,
    namely $\mathcal{D} \in \mathcal{F}_{\alpha}$ for all $\alpha \in \mathcal{A}$.
    Then $\mathcal{G} \coloneqq \mathcal{G}_{\mathcal{D}} = \bigcap_{\alpha \in \mathcal{A}}\mathcal{F}_{\alpha}$ is the smallest sigma-field of subsets of $\mathcal{X}$ containing $\mathcal{D}$,
    also known as the $\sigma$-field generated by $\mathcal{D}$.
\end{definition}

\subsubsection{The canonical probability space}
In the uncountable case of $\Omega = [0, 1]$,
we can start with the collection $\mathcal{D}$ of all intervals $(a, b]$ with $0 \leq a < b \leq 1$ and to define $\P((a, b]) \coloneqq b - a$,
the length of $(a, b]$.
By \autoref{prob:def:siggenfldD} one can uniquely define the sigma-field generated by $\mathcal{D}$,
it is known as the Borel $\sigma$-field $\mathcal{B}[0, 1]$.












\end{document}